\section{Every moment and the exact packing dimension}\label{sec:moments}

The lower bound of Section~\ref{sec:lower} gives only a polynomially small Abel-averaged mass near the critical distance. For the almost Mathieu operator, along sufficiently large scales with $b_n\ge c>h>\log\lambda$, this section proves that the mass beyond $Te^{-hq_n}$ tends to one uniformly in phase and in $e^{hq_n}\le T\le|\delta_n|^{-1}$. Two facts drive the proof. First, the rational discriminant is asymptotically arcsine distributed under the single-site spectral measure, so most of that measure lies away from the parabolic band edges (Lemma~\ref{lem:arcsine}). Second, at such energies the square-summable half-line solution of the equation with energy $E+i/T$---the Weyl solution---remains spread over a long run of elliptic blocks (Lemma~\ref{lem:run}). These facts give Theorem~\ref{thm:fraction}, Theorem~\ref{thm:B} and the lower moment bounds. A diagonal resolvent estimate from the same runs gives the lower local-dimension bound; the upper transport bound gives the upper packing-dimension bound. The dependencies are: Lemmas~\ref{lem:arcsine} and~\ref{lem:run} give Theorem~\ref{thm:fraction}, which gives Corollary~\ref{cor:moments} and Theorems~\ref{thm:A} and~\ref{thm:B}; Lemma~\ref{lem:run} gives Lemma~\ref{lem:localmass} and part~1 of Theorem~\ref{thm:C}; Theorem~\ref{thm:A} gives Proposition~\ref{prop:packupper} and part~2 of Theorem~\ref{thm:C}.

Throughout, $\lambda>1$, $v(x)=2\lambda\cos2\pi x$ and $L=\log\lambda$, except in Subsection~\ref{sec:run}, where $v$ is any Lipschitz function. For a phase $\theta$ we write $\mu^\theta$ for the spectral measure of $\delta_0$ for $H_{\lambda,\alpha,\theta}$, and
\[
G(k;z)=\langle\delta_k,(H_{\lambda,\alpha,\theta}-z)^{-1}\delta_0\rangle\qquad(k\in\Z,\ \operatorname{Im}z>0).
\]
Two identities connect these objects with transport. The first is Parseval's formula, already used in Section~\ref{sec:lower}:
\begin{equation}\label{eq:parseval}
a(k,T)=\frac1{\pi T}\int_\R|G(k;E+i/T)|^2\,dE .
\end{equation}
It follows from $(H-E-i/T)^{-1}=i\int_0^\infty e^{-t/T}e^{itE}e^{-itH}\,dt$ and Plancherel's theorem. The second is the spectral theorem: for $z=E+i\eps$ with $\eps>0$,
\begin{equation}\label{eq:poisson}
\sum_{k\in\Z}|G(k;z)|^2=\frac{\operatorname{Im}G(0;z)}{\eps},\qquad
\frac1\pi\operatorname{Im}G(0;z)=\int_\R P_\eps(E-E')\,d\mu^\theta(E'),\qquad P_\eps(x)=\frac{\eps}{\pi(x^2+\eps^2)}.
\end{equation}

\subsection{The arcsine law}

Fix a scale $n$ and write $q=q_n$, $p=p_n$, $\delta=\delta_n$. For the phase $\theta$ put
\[
\Delta_n(E)=P_{q}(E)-2\lambda^{q}\cos(2\pi q\theta),
\]
the discriminant of the periodic operator $H_{\lambda,p/q,\theta}$ (Lemma~\ref{lem:chambers}). For that periodic operator the distribution of $\Delta_n$ under the spectral measure of $\delta_0$ is exactly the arcsine law: by~\eqref{eq:ratmagic}, $\Delta_n(H_{\lambda,p/q,\theta})=S^q+S^{-q}$. The next lemma transfers this to the irrational operator at every scale where $|\delta_n|$ is much smaller than $\lambda^{-q_n}$.

Let $\sigma$ be the arcsine law on $[-2,2]$, with density $1/(\pi\sqrt{4-x^2})$. Its moments are $m_k=\binom{k}{k/2}$ for even $k$ and $m_k=0$ for odd $k$, and $\sigma([-a,a])=\frac2\pi\arcsin(a/2)$ for $0\le a\le2$.

\begin{lemma}[Arcsine law]\label{lem:arcsine}
Let $c>L$ and let $\mathcal N$ be an infinite set of indices with $b_n\ge c$ for $n\in\mathcal N$. For every phase $\theta$ let $\nu_n$ be the image of $\mu^\theta$ under $\Delta_n$. Then $\int x^k\,d\nu_n\to m_k$ for every integer $k\ge0$, uniformly in $\theta$, as $n\to\infty$ in $\mathcal N$. Consequently $\nu_n$ converges weakly to $\sigma$, and for every $0\le\eta<2$,
\[
\lim_{n\in\mathcal N}\mu^\theta\bigl(\{E:|\Delta_n(E)|\le2-\eta\}\bigr)=\frac2\pi\arcsin\Bigl(1-\frac\eta2\Bigr).
\]
\end{lemma}

\begin{proof}
Fix $h\in(L,c)$. By Proposition~\ref{prop:poly}, for large $n$ and every phase,
\[
\Delta_n(H)=Q-V(0)=D+F+R,\qquad D=V-V(0),\quad F=S^q+S^{-q},
\]
where $D$ is diagonal and $R$ has range at most $q$ and norm at most $C_hq^{5/2}e^{hq}|\delta|$. Each of $D$, $F$ and $R$ maps a vector supported in $[-jq,jq]$ to a vector supported in $[-(j+1)q,(j+1)q]$.

Fix $k\ge1$. In the expansion of $\Delta_n(H)^k\delta_0$ the operator $D$ only acts on vectors supported in $[-kq,kq]$, where
\[
|D(i)|=2\lambda^q\bigl|\cos(2\pi(q\theta+i\delta))-\cos(2\pi q\theta)\bigr|\le4\pi\lambda^q\,kq|\delta|.
\]
So $\Delta_n(H)^k\delta_0=(D_k+F+R)^k\delta_0$ with $D_k=D\un_{[-kq,kq]}$. Put $\eps_n=4\pi kq\lambda^q|\delta|+C_hq^{5/2}e^{hq}|\delta|$, so that $\|D_k+R\|\le\eps_n$, and recall $\|F\|\le2$. Expanding the product,
\[
\bigl\|\Delta_n(H)^k\delta_0-F^k\delta_0\bigr\|\le(2+\eps_n)^k-2^k .
\]
Since $b_n\ge c>h>L$, both $q\lambda^q|\delta|=qe^{(L-b_n)q}$ and $q^{5/2}e^{hq}|\delta|=q^{5/2}e^{(h-b_n)q}$ tend to $0$, uniformly in $\theta$, so $\eps_n\to0$. Next, $F^k\delta_0=\sum_i\binom{k}{(k+i)/2}\delta_{iq}$, the sum running over $|i|\le k$ with $i\equiv k\pmod 2$, so $\langle\delta_0,F^k\delta_0\rangle=m_k$. Since $\int x^k\,d\nu_n=\langle\delta_0,\Delta_n(H)^k\delta_0\rangle$, we obtain $|\int x^k\,d\nu_n-m_k|\le(2+\eps_n)^k-2^k\to0$.

The law $\sigma$ has compact support, so it is determined by its moments~\cite[Theorem~30.1]{Billingsley1995}, and convergence of all moments implies $\nu_n\Rightarrow\sigma$~\cite[Theorem~30.2]{Billingsley1995}. Since $\sigma$ has no atoms, weak convergence gives $\nu_n(J)\to\sigma(J)$ for every interval $J$. Take $J=[-2+\eta,2-\eta]$. The convergence is uniform in $\theta$: otherwise there would be $\eps_0>0$, indices $n_k\to\infty$ in $\mathcal N$ and phases $\theta_k$ with $|\nu^{\theta_k}_{n_k}(J)-\sigma(J)|\ge\eps_0$; but the moments of $\nu^{\theta_k}_{n_k}$ still converge to $m_j$ for every $j$, so $\nu^{\theta_k}_{n_k}\Rightarrow\sigma$ and $\nu^{\theta_k}_{n_k}(J)\to\sigma(J)$, a contradiction.
\end{proof}

\begin{lemma}[Uniform arcsine bound]\label{lem:arcsineunif}
Let $c$ and $\mathcal N$ be as in Lemma~\ref{lem:arcsine}, and let $0\le\eta<2$. Then
\[
\liminf_{n\in\mathcal N}\ \inf_{\theta\in\T}\ \mu^\theta\bigl(\{E:|\Delta_n(E)|\le2-\eta\}\bigr)\ge\frac2\pi\arcsin\Bigl(1-\frac\eta2\Bigr).
\]
\end{lemma}

\begin{proof}
Put $a=2-\eta$ and fix $\tau\in(0,a)$. Let $f$ be the continuous function that equals $1$ on $[-a+\tau,a-\tau]$, vanishes outside $(-a,a)$ and is linear on the two remaining intervals, so $0\le f\le\un_{[-a,a]}$. By the Weierstrass theorem there is a polynomial $p_0$ with $|p_0-f|\le\tau$ on $[-3,3]$. Let $d$ be its degree and $A$ a constant with $|p_0(x)|\le A(|x|/3)^d$ for $|x|\ge3$. For an integer $K$ with $2K\ge d$ put
\[
p(x)=p_0(x)-\tau-A(x/3)^{2K}.
\]
Then $p\le\un_{[-a,a]}$ on $\R$: on $[-3,3]$ we have $p\le p_0-\tau\le f\le\un_{[-a,a]}$, and for $|x|\ge3$ we have $p(x)\le|p_0(x)|-A(|x|/3)^{2K}\le0$. Hence, with $\nu_n=\nu_n^\theta$ the image of $\mu^\theta$ under $\Delta_n$,
\[
\mu^\theta\bigl(\{|\Delta_n|\le a\}\bigr)=\int\un_{[-a,a]}\,d\nu_n\ge\int p\,d\nu_n .
\]
The right side is a fixed finite combination of the moments $\int x^k\,d\nu_n$, which converge to $m_k$ uniformly in $\theta$ by the proof of Lemma~\ref{lem:arcsine}. So $\liminf_n\inf_\theta\mu^\theta(\{|\Delta_n|\le a\})\ge\int p\,d\sigma\ge\int f\,d\sigma-2\tau-A(2/3)^{2K}\ge\sigma([-a+\tau,a-\tau])-2\tau-A(2/3)^{2K}$, since $\sigma$ lives on $[-2,2]$. Let $K\to\infty$ ($A$ does not depend on $K$), then $\tau\to0$, and use $\sigma([-a,a])=\frac2\pi\arcsin(a/2)$.
\end{proof}


The proof uses no condition on the phase. The only arithmetic input is $b_n\ge c>L$, which makes the residual $\eps_n$ small on the sites that $k$ steps from the origin can reach.

\subsection{One elliptic run}\label{sec:run}

In this subsection $v:\T\to\R$ is Lipschitz with constant $\mathrm{Lip}(v)$, and $T(E,x)$ and $A_r(E,x)$ are defined with this $v$ as in Section~\ref{sec:lower}. Fix a phase $\theta$ and a convergent $p/q$ of $\alpha$, and put $\delta=q\alpha-p$, so that $x+q\alpha=x+\delta$ in $\T$. Let $z=E+i\eps$ with $E\in\R$ and $\eps>0$, and let $g\in\ell^2(\Z_{\ge0})$ satisfy
\begin{equation}\label{eq:halfline}
g(k+1)+g(k-1)+v(\theta+k\alpha)\,g(k)=z\,g(k)\qquad(k\ge1).
\end{equation}
Cut the half-line into blocks of $q$ sites. Put $x_m=\theta+(mq+1)\alpha$, $\Psi_m=(g(mq+1),g(mq))^{\mathsf T}$ and
\[
\mathcal M_m=\sum_{r=0}^{q-1}|g(mq+1+r)|^2 .
\]
Then $(g(mq+1+r),g(mq+r))^{\mathsf T}=A_r(z,x_m)\Psi_m$ for $0\le r\le q$, and in particular $\Psi_{m+1}=A_q(z,x_m)\Psi_m$. Consecutive blocks see almost the same potential: $x_{m+1}=x_m+\delta$, so the potentials at corresponding sites differ by at most $\mathrm{Lip}(v)|\delta|$. We use the notation $T_r$, $C_r$, $N_\kappa$, $J_0$ and $R_\varrho$ of Lemma~\ref{lem:blocks}, and $d$ for the hyperbolic distance in the upper half-plane, normalized as there.

The first lemma is the telescoping identity behind Lemma~\ref{lem:blocks}, for site-dependent shifts.

\begin{lemma}[Insertion]\label{lem:insert}
Let $T_r$ and $C_r$ be as in Lemma~\ref{lem:blocks} for energies $E-v_r$, $0\le r<q$, let $B\in SL(2,\R)$, and put $u_r=e_1^{\mathsf T}C_rB$. Let $w_0,\dots,w_{q-1}\in\C$ and let $C_r'$ be the products for the energies $E-v_r+w_r$. Then for $0\le r\le q$,
\[
C_r'B=C_rB\,\Pi_r,\qquad\Pi_r=(I+w_{r-1}Q_{r-1})\cdots(I+w_0Q_0),\qquad Q_j=-J_0u_j^{\mathsf T}u_j,
\]
with $\det\Pi_r=1$ and $\|\Pi_r-I\|\le\exp\bigl(\sum_{j<r}|w_j|\,\|u_j\|^2\bigr)-1$.
\end{lemma}

\begin{proof}
The new one-step matrices are $T_j(I+w_jN_\kappa)$. If $C_j'B=C_jB\Pi_j$, then
\[
C_{j+1}'B=T_j(I+w_jN_\kappa)C_jB\Pi_j=C_{j+1}B\,(I+w_jB^{-1}C_j^{-1}N_\kappa C_jB)\,\Pi_j ,
\]
which gives the product formula by induction. The computation in the proof of Lemma~\ref{lem:blocks}, which uses only $\det B=\det C_j=1$, gives $B^{-1}C_j^{-1}N_\kappa C_jB=-J_0u_j^{\mathsf T}u_j$. This matrix has norm $\|u_j\|^2$, and it has trace $-u_jJ_0u_j^{\mathsf T}=0$ and rank one, so $\det(I+w_jQ_j)=1$. Expanding the product gives the norm bound.
\end{proof}

\begin{lemma}[Perturbed rotations]\label{lem:rot}
Let $|\sin\varrho|\ge s>0$ and let $M=R_\varrho(I+Y)\in SL(2,\R)$ with $\|Y\|\le y\le s/8$. Then $M$ is elliptic, and its fixed point $w$ in the upper half-plane satisfies $d(i,w)\le12y/s$.
\end{lemma}

\begin{proof}
The unitary matrix $\mathcal C=2^{-1/2}\begin{pmatrix}1&-i\\1&i\end{pmatrix}$ acts as the Cayley map from the upper half-plane to the unit disc, sends $i$ to $0$, and is an isometry for the hyperbolic metrics. It conjugates $R_\varrho$ to $\mathrm{diag}(e^{-i\varrho},e^{i\varrho})$ and $M$ to a matrix $\begin{pmatrix}a&b\\\bar b&\bar a\end{pmatrix}$ with $|a|^2-|b|^2=1$. Since $\mathcal C$ is unitary, $\|\mathcal CM\mathcal C^{-1}-\mathcal CR_\varrho\mathcal C^{-1}\|=\|R_\varrho Y\|\le y$, so $|b|\le y$ and $|\operatorname{Im}a|\ge s-y\ge\frac78s>|b|$. Hence $M$ is elliptic. The fixed point of $\zeta\mapsto(a\zeta+b)/(\bar b\zeta+\bar a)$ in the disc solves $\bar b\zeta^2-2i(\operatorname{Im}a)\zeta-b=0$ and has modulus
\[
t=\frac{|b|}{|\operatorname{Im}a|+\sqrt{(\operatorname{Im}a)^2-|b|^2}}\le\frac{8y}{7s}\le\frac14 .
\]
Its hyperbolic distance from $0$ is $\log\frac{1+t}{1-t}\le\frac{2t}{1-t}\le\frac83t\le12y/s$.
\end{proof}

\begin{lemma}[Pair inequality]\label{lem:pair}
Let $G$ be a real symmetric positive semidefinite $2\times2$ matrix with trace $S$, and let $R_\varrho$ be a rotation. Then $\Phi^*(G+R_\varrho^{\mathsf T}GR_\varrho)\Phi\ge S(1-|\cos\varrho|)\|\Phi\|^2$ for every $\Phi\in\C^2$.
\end{lemma}

\begin{proof}
Write $G=\frac S2I+gF_\phi$ with $F_\phi=\begin{pmatrix}\cos\phi&\sin\phi\\\sin\phi&-\cos\phi\end{pmatrix}$ and $0\le g\le S/2$. Since $R_\varrho^{\mathsf T}F_\phi R_\varrho=F_{\phi-2\varrho}$ and $F_\phi+F_{\phi-2\varrho}=2\cos\varrho\,F_{\phi-\varrho}$, the real symmetric matrix $G+R_\varrho^{\mathsf T}GR_\varrho=SI+2g\cos\varrho\,F_{\phi-\varrho}$ has eigenvalues $S\pm2g|\cos\varrho|\ge S(1-|\cos\varrho|)$. A real symmetric matrix bounds its Hermitian form on $\C^2$ below by its smallest eigenvalue.
\end{proof}

The point of Lemma~\ref{lem:pair} is that it holds for every polarization $\Phi$. A single block may carry little mass, for instance when the solution is a standing wave with a node there, but two adjacent blocks cannot both do so.

\begin{lemma}[Elliptic propagation]\label{lem:run}
Let $0<\eta\le1$ and assume $|\tr A_q(E,x_0)|\le2-\eta$. Let $B_0\in SL(2,\R)$ map $i$ to the fixed point of $A_q(E,x_0)$ in the upper half-plane, and put
\[
S_0=\sum_{r=0}^{q-1}\|e_1^{\mathsf T}A_r(E,x_0)B_0\|^2,\qquad\Phi_0=B_0^{-1}\Psi_0,
\]
\[
\xi=(\mathrm{Lip}(v)+1)\max(\eps,|\delta|),\qquad m_*=\Bigl\lfloor\frac{\eta^{3/2}}{2^{10}\xi S_0}\Bigr\rfloor .
\]
If $m_*\ge8$, then:
\begin{enumerate}
\item for $0\le m<m_*$, $\ \mathcal M_m+\mathcal M_{m+1}\ge2^{-8}\eta\,S_0\|\Phi_0\|^2$ and $\mathcal M_m\le16S_0\|\Phi_0\|^2$;
\item with $m_0=\lceil m_*/2\rceil$,
\[
\sum_{k>m_0q}|g(k)|^2\ge\frac{\eta^{5/2}}{2^{22}}\cdot\frac{\|\Phi_0\|^2}{\xi};
\]
\item $\sum_{k>m_*q}|g(k)|^2\le2\|\Phi_0\|^2/\eps$, $\ \sum_{k\ge1}|g(k)|^2\le(1/\xi+2/\eps)\|\Phi_0\|^2$, and $|g(0)|\le\|B_0\|\,\|\Phi_0\|$.
\end{enumerate}
\end{lemma}

By Lemma~\ref{lem:blocks}(1), $S_0=|\partial_E\tr A_q(E,x_0)|/|\sin\varrho_0|$, where $\varrho_0$ is the rotation angle of $A_q(E,x_0)$. For fixed $\eta$, the run consists of $m_*$ blocks and has length comparable to $q/(\xi S_0)$, with comparison constants depending on $\eta$. Its length is set by the larger of the damping $\eps$ and the drift $|\delta|$; it is not a fixed function of the scale. When $\eps\ge|\delta|$, the total mass in part~3 is at most $3(\mathrm{Lip}(v)+1)\|\Phi_0\|^2/\xi$, so part~2 says that the far half of the run carries a fixed fraction of the mass of $g$ on $\Z_{>0}$. No arithmetic condition on $\theta$ enters, and nothing is assumed about the solution beyond the run.

\begin{proof}
\emph{Frames.} Write $A_m=A_q(E,x_m)$ for the real-energy block monodromies. Suppose $A_m$ is elliptic and $B_m\in SL(2,\R)$ maps $i$ to its fixed point, so that $B_m^{-1}A_mB_m=R_{\varrho_m}$ by Lemma~\ref{lem:blocks}(4). Put
\[
u_r^{(m)}=e_1^{\mathsf T}A_r(E,x_m)B_m,\qquad G_m=\sum_{r=0}^{q-1}(u_r^{(m)})^{\mathsf T}u_r^{(m)},\qquad S_m=\tr G_m,\qquad\Phi_m=B_m^{-1}\Psi_m .
\]

\emph{One step.} Apply Lemma~\ref{lem:insert} with $B=B_m$ and the shifts $w_r=v(x_m+r\alpha)-v(x_{m+1}+r\alpha)$, which satisfy $|w_r|\le\mathrm{Lip}(v)|\delta|$. It gives, for $0\le r\le q$,
\[
A_r(E,x_{m+1})B_m=A_r(E,x_m)B_m\,\Pi^{(m)}_r,\qquad\|\Pi^{(m)}_r-I\|\le y_m:=e^{\mathrm{Lip}(v)|\delta|S_m}-1 .
\]
In particular $B_m^{-1}A_{m+1}B_m=R_{\varrho_m}\Pi^{(m)}_q$ and $|\tr A_{m+1}-\tr A_m|\le2y_m$. Put $s_0=\sqrt\eta/2$. If $|\sin\varrho_m|\ge s_0$ and $y_m\le s_0/8$, Lemma~\ref{lem:rot} shows that $A_{m+1}$ is elliptic with fixed point $B_m(w_m)$, where $d(i,w_m)\le d_m:=12y_m/s_0$. Choose $N_m\in SL(2,\R)$ with $N_m(i)=w_m$; by Lemma~\ref{lem:blocks}(4), $\|N_m\|^2=\|N_m^{-1}\|^2=e^{d(i,w_m)}\le e^{d_m}$. Then $B_{m+1}=B_mN_m$ maps $i$ to the fixed point of $A_{m+1}$.

\emph{The damped step.} Lemma~\ref{lem:insert} with $B=B_m$ and the constant shift $w_r=i\eps$ gives $A_q(z,x_m)B_m=B_mR_{\varrho_m}\Pi'_m$ with $\det\Pi'_m=1$ and $\|\Pi'_m-I\|\le e^{\eps S_m}-1$. Hence
\[
\Phi_{m+1}=N_m^{-1}R_{\varrho_m}\Pi'_m\,\Phi_m .
\]

\emph{Bootstrap.} We show by induction that for $0\le m\le m_*$ the frame $B_m$ is defined and
\[
(a_m)\qquad|\tr A_m|\le2-\eta/2,\qquad S_0/2\le S_m\le2S_0,\qquad\|\Phi_0\|/2\le\|\Phi_m\|\le2\|\Phi_0\|.
\]
$(a_0)$ holds. Assume $(a_j)$ for all $j<m$, where $m\le m_*$. Note that $\eps\le\xi$ and $\mathrm{Lip}(v)|\delta|\le\xi$, and that $m_*\ge1$ forces $\xi S_0\le2^{-10}\eta^{3/2}$. For $j<m$ we have $|\cos\varrho_j|\le1-\eta/4$, hence $\sin^2\varrho_j\ge\eta/4$ and $|\sin\varrho_j|\ge s_0$. Also $\mathrm{Lip}(v)|\delta|S_j\le2\xi S_0\le1$, and $e^x-1\le2x$ for $0\le x\le1$, so
\[
y_j\le4\xi S_0\le2^{-8}\eta^{3/2}\le s_0/8 .
\]
So the frames $B_1,\dots,B_m$ are defined. Summing the one-step estimates over $j<m$:
\begin{itemize}
\item $|\tr A_m-\tr A_0|\le2\sum_{j<m}y_j\le8m_*\xi S_0\le2^{-7}\eta^{3/2}\le\eta/2$, so $|\tr A_m|\le2-\eta/2$.
\item $u_r^{(j+1)}=u_r^{(j)}\Pi^{(j)}_rN_j$, so, since $\|u\Pi^{(j)}_r\|\ge\|u\|-\|u(\Pi^{(j)}_r-I)\|$ and $\|uN_j\|\ge\|u\|/\|N_j^{-1}\|$, $S_{j+1}/S_j\in[(1-y_j)^2e^{-d_j},(1+y_j)^2e^{d_j}]$ and $|\log(S_{j+1}/S_j)|\le3y_j+d_j\le27y_j/\sqrt\eta$. Hence $|\log(S_m/S_0)|\le108\,m_*\xi S_0/\sqrt\eta\le108\eta/2^{10}<\log2$.
\item A $2\times2$ matrix of determinant one has $\|M^{-1}\|=\|M\|$, and $\|\Pi'_j\|\le e^{\eps S_j}$. So $|\log(\|\Phi_{j+1}\|/\|\Phi_j\|)|\le\eps S_j+d_j/2\le2\xi S_0+12y_j/\sqrt\eta\le50\xi S_0/\sqrt\eta$, and $|\log(\|\Phi_m\|/\|\Phi_0\|)|\le50\eta/2^{10}<\log2$.
\end{itemize}
This proves $(a_m)$.

\emph{Pairs of blocks.} Fix $0\le m<m_*$ and put $\eps_2=e^{3\xi S_m}-1$. Since $3\xi S_m\le6\xi S_0\le1$, we have $\eps_2\le12\xi S_0\le12\cdot2^{-10}\eta^{3/2}$. By Lemma~\ref{lem:insert} with the shift $i\eps$,
\[
g(mq+1+r)=e_1^{\mathsf T}A_r(z,x_m)B_m\Phi_m=u_r^{(m)}(\Phi_m+e_r),\qquad\|e_r\|\le\eps_2\|\Phi_m\| .
\]
For the next block, $\Psi_{m+1}=B_mR_{\varrho_m}\Pi'_m\Phi_m$. Apply Lemma~\ref{lem:insert} first with the phase shifts, which gives $A_r(E,x_{m+1})B_m=A_r(E,x_m)B_m\Pi^{(m)}_r$, and then with the shift $i\eps$, for which the rows $u_r^{(m)}\Pi^{(m)}_r$ have squared norms at most $(1+y_m)^2\|u_r^{(m)}\|^2\le2\|u_r^{(m)}\|^2$. This gives
\[
g((m+1)q+1+r)=u_r^{(m)}\,\Pi^{(m)}_r\Pi''_r\,R_{\varrho_m}\Pi'_m\,\Phi_m=u_r^{(m)}(R_{\varrho_m}\Phi_m+e_r'),
\]
where $\|\Pi''_r-I\|\le e^{2\eps S_m}-1$. The three factors $\Pi^{(m)}_r$, $\Pi''_r$ and $\Pi'_m$ differ from $I$ by at most $e^{\mathrm{Lip}(v)|\delta|S_m}-1$, $e^{2\eps S_m}-1$ and $e^{\eps S_m}-1$, and $R_{\varrho_m}$ is orthogonal. Hence $\|\Pi^{(m)}_r\Pi''_rR_{\varrho_m}\Pi'_m-R_{\varrho_m}\|\le e^{(\mathrm{Lip}(v)|\delta|+3\eps)S_m}-1\le\eps_2$, because $\mathrm{Lip}(v)|\delta|+3\eps\le3\xi$. So $\|e_r'\|\le\eps_2\|\Phi_m\|$.

For every $e\in\C^2$, $\sum_r|u_r^{(m)}e|^2=e^*G_me\le S_m\|e\|^2$, and $\Phi^*G_m\Phi=\sum_r|u_r^{(m)}\Phi|^2$. Apply $\|a+b\|^2\ge\frac12\|a\|^2-\|b\|^2$ to vectors indexed by $0\le r<q$:
\[
\mathcal M_m\ge\tfrac12\Phi_m^*G_m\Phi_m-S_m\eps_2^2\|\Phi_m\|^2,\qquad
\mathcal M_{m+1}\ge\tfrac12\Phi_m^*R_{\varrho_m}^{\mathsf T}G_mR_{\varrho_m}\Phi_m-S_m\eps_2^2\|\Phi_m\|^2 .
\]
By $(a_m)$, $|\cos\varrho_m|\le1-\eta/4$, so Lemma~\ref{lem:pair} and $2\eps_2^2\le\eta/16$ give
\[
\mathcal M_m+\mathcal M_{m+1}\ge\Bigl(\frac\eta8-2\eps_2^2\Bigr)S_m\|\Phi_m\|^2\ge\frac\eta{16}S_m\|\Phi_m\|^2\ge\frac{\eta}{2^7}S_0\|\Phi_0\|^2 .
\]
Also $\mathcal M_m\le(1+\eps_2)^2S_m\|\Phi_m\|^2\le16S_0\|\Phi_0\|^2$. This proves part~1.

\emph{Flux.} Multiply~\eqref{eq:halfline} by $\overline{g(k)}$, sum over $k>n$ for some $n\ge0$, and take imaginary parts. Since $g\in\ell^2$ and $\sum_{k>n}2\operatorname{Re}\bigl(g(k+1)\overline{g(k)}\bigr)$ is real,
\[
\eps\sum_{k>n}|g(k)|^2=\operatorname{Im}\bigl(g(n)\overline{g(n+1)}\bigr).
\]
If $(g(n+1),g(n))^{\mathsf T}=B\Phi$ with $B$ real and $\det B=1$, then $2|\operatorname{Im}(g(n)\overline{g(n+1)})|=|\det(B\Phi,B\bar\Phi)|=|\det(\Phi,\bar\Phi)|\le\|\Phi\|^2$. At $n=m_*q$ this gives $\sum_{k>m_*q}|g(k)|^2\le\|\Phi_{m_*}\|^2/(2\eps)\le2\|\Phi_0\|^2/\eps$.

\emph{Counting.} Let $X=\eta^{3/2}/(2^{10}\xi S_0)$, so that $X\ge m_*\ge8$ and $m_*\ge X/2$. The blocks $m_0,\dots,m_*$ lie in $\{k>m_0q\}$ and contain at least $\lfloor m_*/2\rfloor/2\ge m_*/8$ disjoint pairs $(m,m+1)$ with $m<m_*$. By part~1,
\[
\sum_{k>m_0q}|g(k)|^2\ge\frac{m_*}8\cdot\frac{\eta}{2^8}S_0\|\Phi_0\|^2\ge\frac{X}{16}\cdot\frac{\eta}{2^8}S_0\|\Phi_0\|^2=\frac{\eta^{5/2}}{2^{22}}\cdot\frac{\|\Phi_0\|^2}{\xi}.
\]
For part~3, $\sum_{k\ge1}|g(k)|^2=\sum_{m<m_*}\mathcal M_m+\sum_{k>m_*q}|g(k)|^2$, and $16m_*S_0\le16XS_0=2^{-6}\eta^{3/2}/\xi\le1/\xi$. Finally $g(0)=e_2^{\mathsf T}\Psi_0=e_2^{\mathsf T}B_0\Phi_0$.
\end{proof}

\subsection{Almost all of the packet at resonant times}

From now on $v(x)=2\lambda\cos2\pi x$ again, so $\mathrm{Lip}(v)=4\pi\lambda$. We first compare the traces of the actual $q$-blocks with the rational discriminant.

\begin{lemma}[Block traces]\label{lem:blocktrace}
For every $h>L$ there is $C$ such that for all large $n$, all real $E\in U_h$ and all $x\in\T$,
\[
\bigl|\tr A_{q_n}(E,x)-P_{q_n}(E)+2\lambda^{q_n}\cos(2\pi q_nx)\bigr|\le C\,q_ne^{hq_n}|\delta_n| .
\]
\end{lemma}

\begin{proof}
By Lemma~\ref{lem:chambers}, $P_q(E)-2\lambda^q\cos(2\pi qx)=\tr A^{(p/q)}_q(E,x)$. As in Step~1 of the proof of Proposition~\ref{prop:poly}, write each rational factor as the irrational one plus an error of norm at most $d=4\pi\lambda|\delta|$ and expand. The difference $A^{(p/q)}_q(E,x)-A_q(E,x)$ is the sum of the terms with at least one error, so its norm is at most
\[
\sum_{k=1}^q\binom qkC_h^{k+1}e^{h(q-k)}d^k=C_he^{hq}\bigl[(1+C_hde^{-h})^q-1\bigr]\le C_h^2qd\,e^{hq}e^{C_hqd},
\]
and $qd\le4\pi\lambda$. The trace of a $2\times2$ matrix is at most twice its norm.
\end{proof}

\begin{theorem}[Almost all of the packet moves out]\label{thm:fraction}
Let $\lambda>1$, let $\alpha$ be irrational and let $\theta$ be any phase. Fix $h>L$, $c>h$ and $0<\eta\le1/4$. There are $C$ and $n_0$, depending only on $\lambda$, $\alpha$, $h$, $c$ and $\eta$, such that for every $n\ge n_0$ with $b_n\ge c$ and every $T$ with $e^{hq_n}\le T\le|\delta_n|^{-1}$,
\[
P\bigl(Te^{-hq_n},T\bigr)\ge\bigl(1-Cq_ne^{-(h-L)q_n/2}\bigr)\,\mu^\theta\bigl(\{E:|\Delta_n(E)|\le2-3\eta\}\bigr).
\]
\end{theorem}

If $b_n\ge c$ for infinitely many $n$, then along those $n$ the factor $\mu^\theta(\cdot)$ is at least $\frac2\pi\arcsin(1-\frac{3\eta}2)-o(1)$, uniformly in $\theta$, by Lemma~\ref{lem:arcsineunif}; since $\eta$ is arbitrary, the Abel-averaged probability beyond $Te^{-hq_n}$ tends to $1$ (Theorem~\ref{thm:B}). At the resonant time $T=|\delta_n|^{-1}=e^{b_nq_n}$ the distance is $e^{(b_n-h)q_n}=T^{1-h/b_n}$. The mechanism is visible in the proof: at an energy whose first $q$-blocks on both sides are elliptic, the resolvent $G(\cdot;E+i/T)$ has, on each side, comparable mass on every pair of adjacent blocks among at least $c'T/(qe^{h_0q})$ consecutive $q$-blocks, and at most $16S_0\|\Phi_0\|^2$ on each of them (Lemma~\ref{lem:run}(1)); so the at most $Te^{-hq}/q+1$ blocks within distance $Te^{-hq}$ carry only a fraction $O(qe^{-(h-h_0)q})$ of it, with $h_0=(L+h)/2$. The mass beyond the run, which the earlier version of this argument counted against the fraction, is itself beyond $Te^{-hq}$.

\begin{proof}
Let $h_0=(L+h)/2$ and $U=U_{h_0}$. By Lemma~\ref{lem:growth} and Proposition~\ref{prop:poly}(2) there is $C$ with $\|A_r(E,x)\|\le Ce^{h_0r}$ for real $E\in U$, all $x$ and $r\ge0$, and $\sup_U|P_q'|\le Cqe^{h_0q}$. Write $q=q_n$, $\delta=\delta_n$, $b=b_n$ and $\eps=1/T$, so that $|\delta|\le\eps\le e^{-hq}$, and put
\[
N=Te^{-hq},\qquad I=\{E\in U:\ |\Delta_n(E)|\le2-2\eta\}.
\]
Below, $C$ and $c'$ denote positive constants depending only on $\lambda$, $\alpha$, $h$ and $\eta$, and ``for large $n$'' means for $n$ beyond a threshold depending only on $\lambda$, $\alpha$, $h$, $c$ and $\eta$; nothing depends on $\theta$ or $T$. Note that $h_0-h=-(h-L)/2$.

\emph{Step 1: two elliptic first blocks.} Fix $E\in I$ and $z=E+i\eps$. For $k\ge0$ put $g_+(k)=G(k;z)$ and $g_-(k)=G(-k;z)$. Both lie in $\ell^2(\Z_{\ge0})$. The function $g_+$ satisfies~\eqref{eq:halfline} for the phase $\theta$. Since the cosine is even, $v(\theta-k\alpha)=v(-\theta+k\alpha)$, so $g_-$ satisfies~\eqref{eq:halfline} for the phase $-\theta$. The first blocks are $A_q(E,\pm\theta+\alpha)$. Since $q(\pm\theta+\alpha)=\pm q\theta+p+\delta$, Lemma~\ref{lem:blocktrace} and $|\cos(2\pi(\pm q\theta+\delta))-\cos(2\pi q\theta)|\le2\pi|\delta|$ give
\[
\bigl|\tr A_q(E,\pm\theta+\alpha)-\Delta_n(E)\bigr|\le Cqe^{h_0q}|\delta|+4\pi\lambda^q|\delta| .
\]
The right side tends to $0$ because $b\ge c>h_0>L$. So for large $n$ both first blocks satisfy $|\tr A_q|\le2-\eta$.

\emph{Step 2: the run is long and the near blocks are few.} Apply Lemma~\ref{lem:run} to $g_\pm$, with this $\eta$, the damping $\eps$ and $\xi=(4\pi\lambda+1)\eps$ (recall $\eps\ge|\delta|$), and let $B_0^\pm$, $S_0^\pm$, $\Phi_0^\pm$, $m_*^\pm$, $m_0^\pm$ be the corresponding data. From $|\tr A_q|\le2-\eta$ we get $|\sin\varrho_0^\pm|\ge\sqrt\eta/2$, and the product rule gives $|\partial_E\tr A_q(E,x)|\le2qC^2e^{h_0q}$. By Lemma~\ref{lem:blocks}(1), $S_0^\pm\le Cqe^{h_0q}$. Hence
\[
\frac{\eta^{3/2}}{2^{10}\xi S_0^\pm}\ge\frac{c'T}{qe^{h_0q}}\ge\frac{c'e^{(h-h_0)q}}{q},
\]
so for large $n$ we have $m_*^\pm\ge8$, and $m_0^\pm q\ge m_*^\pm q/2\ge c'Te^{-h_0q}/4\ge Te^{-hq}=N$.

\emph{Far mass.} Lemma~\ref{lem:run}(2) and $m_0^\pm q\ge N$ give
\begin{equation}\label{eq:farlower}
\sum_{k>N}|g_\pm(k)|^2\ge\frac{\eta^{5/2}}{2^{22}}\cdot\frac{\|\Phi_0^\pm\|^2}{\xi}.
\end{equation}

\emph{Near mass.} The sites $1\le k\le N$ lie in the blocks $\{mq+1,\dots,mq+q\}$ with $0\le m\le(N-1)/q$, which are at most $N/q+1$ blocks, and $(N-1)/q<m_0^\pm\le m_*^\pm-1$ because $m_*^\pm\ge8$. So Lemma~\ref{lem:run}(1) applies to each of them, and with~\eqref{eq:farlower},
\[
\sum_{k=1}^N|g_\pm(k)|^2\le\Bigl(\frac Nq+1\Bigr)16S_0^\pm\|\Phi_0^\pm\|^2\le2^{26}\eta^{-5/2}\Bigl(\frac Nq+1\Bigr)S_0^\pm\xi\sum_{k>N}|g_\pm(k)|^2 .
\]
Here, using $S_0^\pm\le Cqe^{h_0q}$, $\xi=(4\pi\lambda+1)/T$, $N=Te^{-hq}$ and $T\ge e^{hq}$,
\[
\Bigl(\frac Nq+1\Bigr)S_0^\pm\xi\le C(4\pi\lambda+1)\Bigl(e^{(h_0-h)q}+\frac{qe^{h_0q}}T\Bigr)\le C(4\pi\lambda+1)(1+q)e^{-(h-L)q/2}.
\]

\emph{The site $0$.} By Lemma~\ref{lem:blocks}(4), $\|B_0^+\|^2=e^{d(i,z_+)}$, where $z_+$ is the fixed point of $\mathcal A=A_q(E,\theta+\alpha)$. The identity $\|\mathcal A\|_{\mathrm{HS}}^2=2+4\sinh^2(d(i,z_+))\sin^2\varrho_0^+$ from the proof of Lemma~\ref{lem:blocks}(4) gives $\|B_0^+\|^2\le1+4\|\mathcal A\|/\sqrt\eta\le Ce^{h_0q}$. With part~3 of Lemma~\ref{lem:run} and~\eqref{eq:farlower},
\[
|G(0;z)|^2\le\|B_0^+\|^2\|\Phi_0^+\|^2\le Ce^{h_0q}\xi\sum_{k>N}|g_+(k)|^2\le C(4\pi\lambda+1)e^{-(h-L)q/2}\sum_{k>N}|g_+(k)|^2 .
\]

\emph{Conclusion of Step 2.} Since $g_+(k)=G(k;z)$ and $g_-(k)=G(-k;z)$ for $k\ge1$, adding the last three displays gives, for large $n$,
\[
\sum_{|k|\le N}|G(k;z)|^2\le\kappa_n\sum_{|k|>N}|G(k;z)|^2,\qquad\kappa_n=Cqe^{-(h-L)q/2},
\]
and therefore
\begin{equation}\label{eq:farfraction}
\sum_{|k|>N}|G(k;z)|^2\ge(1-\kappa_n)\sum_{k\in\Z}|G(k;z)|^2\qquad(E\in I).
\end{equation}

\emph{Step 3: Parseval.} By~\eqref{eq:parseval}, \eqref{eq:farfraction} and~\eqref{eq:poisson},
\[
P(N,T)=\frac1{\pi T}\int_\R\sum_{|k|>N}|G(k;E+i\eps)|^2\,dE\ge(1-\kappa_n)\int_I\frac1\pi\operatorname{Im}G(0;E+i\eps)\,dE ,
\]
and the right side equals $(1-\kappa_n)\int_\R\bigl(\int_IP_\eps(E-E')\,dE\bigr)\,d\mu^\theta(E')$.

\emph{Step 4: Poisson smoothing.} Let $r_*>0$ be the distance from $\Sigma$ to $\R\setminus U$, and $r=\min(r_*/2,\eta/\sup_U|P_q'|)$, so that $r\ge\eta/(Cqe^{h_0q})$ for large $n$. If $E'\in\Sigma$ and $|\Delta_n(E')|\le2-3\eta$, then $[E'-r,E'+r]\subset U$, and $\Delta_n'=P_q'$ gives $|\Delta_n|\le2-2\eta$ there. So $[E'-r,E'+r]\subset I$ and
\[
\int_IP_\eps(E-E')\,dE\ge\frac2\pi\arctan\frac r\eps\ge1-\frac{2\eps}{\pi r} .
\]
Since $\eps/r\le Cqe^{h_0q}/(\eta T)\le Cqe^{-(h-L)q/2}$ and $\mu^\theta$ is supported on $\Sigma$, we obtain $P(N,T)\ge(1-\kappa_n)(1-Cqe^{-(h-L)q/2})\,\mu^\theta(\{|\Delta_n|\le2-3\eta\})$, which is the claim after enlarging $C$.
\end{proof}

\begin{remark}[Where the old constant was lost]\label{rem:oldconstant}
An earlier version compared the far mass~\eqref{eq:farlower} with the whole mass $\sum_{k\ge1}|g_\pm(k)|^2\le3(4\pi\lambda+1)\|\Phi_0^\pm\|^2/\xi$ of Lemma~\ref{lem:run}(3), which produced the factor $\eta^{5/2}2^{-24}/(4\pi\lambda+1)$, and then took $\eta=1/4$ and $\mu^\theta(\{|\Delta_n|\le5/4\})\ge1/4$, which gives $2^{-31}/(1+4\pi\lambda)$; an extra factor $1/2$, taken in the eventual Poisson-smoothing bound, produced $2^{-32}/(1+4\pi\lambda)$. The upper bound of Lemma~\ref{lem:run}(3) counts the mass beyond the run, but that mass is also beyond $N$. The flux identity, the Parseval step and the Poisson smoothing lose only factors $1-o(1)$; the only genuine loss is the spectral mass near the band edges, $1-\frac2\pi\arcsin(1-\frac{3\eta}2)\approx\frac2\pi\sqrt{3\eta}$, which disappears as $\eta\to0$.
\end{remark}

\begin{remark}[An annulus at strongly resonant scales]\label{rem:burst}
Theorem~\ref{thm:fraction} has a matching upper half. Let $\theta\in\Ga^*$ and $\beta(\alpha)<\infty$, and fix $B_0>\beta(\alpha)$, $h>L$, $c>h$, $\eps>0$ and $M>0$. Let $\zeta_1$ be given by Theorem~\ref{thm:window} for $\eta=c-L$, this $\eps$ and $M$, and $\Lambda=\max(1,B_0)$. For large $n$ with $b_n\ge c$ we have $b_n\le B_0$ and, since $\theta\in\Ga^*$, $|\sin(2\pi q_n\theta)|\ge e^{-\zeta_1q_n}$; this threshold in $n$ may depend on the fixed phase. For every $T$ with $e^{hq_n}\le T\le|\delta_n|^{-1}$ put
\[
N_-=Te^{-hq_n}=T\lambda^{-q_n}e^{-(h-L)q_n},\qquad N_+=e^{\eps q_n}\bigl(1+T\lambda^{-q_n}\bigr)\le2e^{\eps q_n}T\lambda^{-q_n}.
\]
Theorem~\ref{thm:window}, whose time range contains $T\le|\delta_n|^{-1}\le e^{B_0q_n}$, gives $P(N_+,T)\le Ce^{-Mq_n}$, and Theorem~\ref{thm:B}, which is deduced from Theorem~\ref{thm:fraction} and Lemma~\ref{lem:arcsineunif} in Section~\ref{sec:everymoment}, gives $P(N_-,T)\ge1-o(1)$ as $n\to\infty$ with $b_n\ge c$, uniformly in $T$. So the annulus $N_-<|j|\le N_+$ carries Abel-averaged mass $1-o(1)$: asymptotically the whole Abel-averaged packet lies at distance $T\lambda^{-q_n}$, up to the factors $e^{(h-L)q_n}$ and $2e^{\eps q_n}$, whose exponents can be fixed as small as we like before $n$ is taken large. In particular, for each fixed $s\in(0,1)$, the $s$-quantile of the Abel-averaged distribution $j\mapsto a(j,T)$ of $|j|$ (the median for $s=1/2$) lies in $[N_-,N_+]$ for all large $n$ with $b_n\ge c$, uniformly in $T$. The window begins at $T=e^{hq_n}=\lambda^{q_n}e^{(h-L)q_n}$, later than $\lambda^{q_n}$ by the factor $e^{(h-L)q_n}$ for the fixed $h>L$, because Theorem~\ref{thm:fraction} requires $T\ge e^{hq_n}$. All of this concerns time averages, not the instantaneous position of the packet.
\end{remark}

\begin{remark}[Numerical illustration; not used in any proof]\label{rem:fractionnumerics}
We propagated $\delta_0$ for $H_{\lambda,\alpha,\theta}$, $\alpha=p/q+\eta_\alpha$, by a Chebyshev expansion of $e^{-itH}$ on a box of half-width at least $6T\lambda^{-q}+300$ (the mass ever seen within $30$ sites of the box edge was below $10^{-150}$), and computed $a(j,T)$ by the trapezoid rule in $t$ on $[0,7T]$ with step $25$ (the neglected Abel weight is $e^{-14}$; the computed total mass was within $1.3\cdot10^{-3}$ of $1$). With $\delta=q\alpha-p$ and $\ell=T\lambda^{-q}$ we found:
\begin{center}
\begin{tabular}{cccccccc}
$\lambda$&$p/q$&$\theta$&$T$&$P(0.01\ell,T)$&$P(0.1\ell,T)$&$P(\ell,T)$&median of $|j|$\\\hline
$1.5$&$2/5$&$0.2137$&$|\delta|^{-1}=2000$&$0.945$&$0.674$&$0.049$&$0.21\,\ell$\\
$1.5$&$2/5$&$0$&$|\delta|^{-1}=2000$&$0.891$&$0.473$&$10^{-8}$&$0.09\,\ell$\\
$1.5$&$2/5$&$0.2137$&$0.2|\delta|^{-1}=400$&$0.911$&$0.686$&$0.115$&$0.27\,\ell$\\
$2$&$2/5$&$0.2137$&$|\delta|^{-1}\approx6667$&$0.951$&$0.743$&$0.002$&$0.26\,\ell$\\
$2.5$&$2/5$&$0.2137$&$|\delta|^{-1}=20000$&$0.954$&$0.747$&$10^{-6}$&$0.27\,\ell$\\
$1.5$&$3/8$&$0.2137$&$|\delta|^{-1}=6250$&$0.869$&$0.498$&$10^{-25}$&$0.10\,\ell$\\
$1.2$&$2/5$&$0.2137$&$|\delta|^{-1}=300\lambda^5$&$0.925$&$0.563$&$10^{-23}$&$0.12\,\ell$\\
$1.2$&$3/8$&$0.2137$&$|\delta|^{-1}=300\lambda^8$&$0.859$&$0.452$&$10^{-11}$&$0.08\,\ell$\\
$1.2$&$5/13$&$0.2137$&$|\delta|^{-1}=300\lambda^{13}$&$0.789$&$0.361$&$10^{-19}$&$0.08\,\ell$\\
$1.2$&$8/21$&$0.2137$&$|\delta|^{-1}=300\lambda^{21}$&$0.927$&$0.698$&$0.114$&$0.24\,\ell$
\end{tabular}
\end{center}
So the Abel-averaged distribution of $|j|$ lives on the scale $T\lambda^{-q}$, with a profile that depends on $\lambda$, $q$ and $\theta$ (for $\lambda=1.2$ the median moved non-monotonically in $q$, between $0.08\,\ell$ and $0.24\,\ell$), and the fraction beyond a fixed multiple $c\,T\lambda^{-q}$ is far from $1$ unless $c$ is small. This is consistent with Theorem~\ref{thm:B}, whose threshold $Te^{-hq_n}$ is smaller than $T\lambda^{-q_n}$ by the factor $e^{-(h-\log\lambda)q_n}\to0$, and indicates that some such vanishing factor is needed for a fraction tending to $1$. At these denominators $e^{-(h-\log\lambda)q/2}$ is not small, so the computation illustrates the theorem but does not test its rate.

\end{remark}

\subsection{Every moment}\label{sec:everymoment}

\begin{corollary}\label{cor:moments}
Let $\lambda>1$ and let $\alpha$ be irrational. For every phase $\theta$ and every $p>0$,
\[
\beta^+(p)\ge\max\Bigl\{0,\ 1-\frac{\log\lambda}{\beta(\alpha)}\Bigr\},
\]
where the right side is read as $0$ when $\beta(\alpha)=0$ and as $1$ when $\beta(\alpha)=\infty$. If $\beta(\alpha)>\log\lambda$, moreover $\limsup_{T\to\infty}P(T^a,T)>0$ for every $0<a<1-\log\lambda/\beta(\alpha)$.
\end{corollary}

\begin{proof}
\emph{Nonnegativity.} The vector $\delta_0$ is not an eigenvector of $H$, since $H\delta_0$ has nonzero entries at $\pm1$. So $\mu^\theta$ is not a point mass and $\sum_E\mu^\theta(\{E\})^2<1$. Now
\[
a(0,T)=\iint\frac{2/T}{2/T+i(E-E')}\,d\mu^\theta(E)\,d\mu^\theta(E'),
\]
and the integrand has modulus at most one and tends to $\un_{E=E'}$, so $a(0,T)\to\sum_E\mu^\theta(\{E\})^2$ by dominated convergence. Hence $M_p(T)\ge P(0,T)=1-a(0,T)$ is bounded below for large $T$, and $\beta^+(p)\ge0$.

\emph{The main bound.} Assume $\beta(\alpha)>L$. Fix $h\in(L,\beta(\alpha))$ and $c\in(h,\beta(\alpha))$. Since $\limsup b_n=\beta(\alpha)$, there are $n_j\to\infty$ with $b_{n_j}\ge c$ and $b_{n_j}\to\beta(\alpha)$. Apply Theorem~\ref{thm:fraction} with $\eta=1/4$ at $T_j=|\delta_{n_j}|^{-1}$, and put $N_j=T_je^{-hq_{n_j}}=e^{(b_{n_j}-h)q_{n_j}}$. By Lemma~\ref{lem:arcsine}, $\mu^\theta(\{|\Delta_{n_j}|\le5/4\})\to\frac2\pi\arcsin\frac58>\frac25$, and the factor $1-Cq_{n_j}e^{-(h-L)q_{n_j}/2}$ tends to $1$, so $P(N_j,T_j)\ge c_0:=1/4$ for large $j$. Hence
\[
M_p(T_j)\ge N_j^p\,P(N_j,T_j)\ge c_0N_j^p,\qquad\frac{\log N_j}{\log T_j}=\frac{b_{n_j}-h}{b_{n_j}}\to1-\frac h{\beta(\alpha)},
\]
with limit $1$ if $\beta(\alpha)=\infty$. So $\beta^+(p)\ge1-h/\beta(\alpha)$, and letting $h\downarrow L$ proves the bound. If $a<1-L/\beta(\alpha)$, choose $h$ with $a<1-h/\beta(\alpha)$; then $T_j^a\le N_j$ for large $j$, so $P(T_j^a,T_j)\ge c_0$.
\end{proof}

Damanik and Tcheremchantsev~\cite{DT2007} also use the exponent $\alpha_l^+$, the supremum of the $a$ for which $P(T^a,T)$ decays at no polynomial rate along some sequence of times; always $\alpha_l^+\le\au$. The second statement of Corollary~\ref{cor:moments} gives $\alpha_l^+\ge1-\log\lambda/\beta(\alpha)$. Thus, for $\theta\in\Ga^*$ and $\beta(\alpha)<\infty$, $\alpha_l^+=\au$. When $\gamma:=1-\log\lambda/\beta(\alpha)>0$, Theorem~\ref{thm:B} in fact gives, for every $0<a<\gamma$, Abel-averaged outside probability tending to one beyond $T^a$ along a sequence of averaging times. This does not assert a positive limsup for $P(T^\gamma,T)$ at the exact critical radius.

\begin{proof}[Proof of Theorem~\ref{thm:A}]
Part~1. Theorem~\ref{thm:upper} gives $P(T^a,T)=O(T^{-m})$ for every $a>\max\{0,1-L/\beta(\alpha)\}$ and every $m>0$, and $\au\le\max\{0,1-L/\beta(\alpha)\}$. Since $\beta^+(p)\le\au$ by~\eqref{eq:exponents}, the same bound holds for every $\beta^+(p)$. Corollary~\ref{cor:moments} gives the reverse inequality $\beta^+(p)\ge\max\{0,1-L/\beta(\alpha)\}$ at every phase, and $\au\ge\beta^+(p)$.

Part~2. Corollary~\ref{cor:moments} gives $\beta^+(p)\ge1$ at every phase, and Lemma~\ref{lem:moments}(1) gives $\beta^+(p)\le1$. Hence $\beta^+(p)=\au=1$. This recovers~\cite[Corollary~3]{JZ2022}.
\end{proof}

\begin{proof}[Proof of Theorem~\ref{thm:B}]
Let $\mathcal N=\{n:b_n\ge c\}$; if it is finite there is nothing to prove. Fix $\tau\in(0,1)$ and choose $\eta\in(0,1/4]$ with $\frac2\pi\arcsin(1-\frac{3\eta}2)\ge1-\tau/2$. By Lemma~\ref{lem:arcsineunif}, for large $n\in\mathcal N$ and every phase, $\mu^\theta(\{|\Delta_n|\le2-3\eta\})\ge1-3\tau/4$. Theorem~\ref{thm:fraction} with this $\eta$ gives, for large $n\in\mathcal N$, every phase and every $T\in[e^{hq_n},|\delta_n|^{-1}]$,
\[
P_\theta\bigl(Te^{-hq_n},T\bigr)\ge\bigl(1-Cq_ne^{-(h-L)q_n/2}\bigr)(1-3\tau/4)\ge1-\tau .
\]
For a packet started at $\delta_j$, the translation $U^j$ of the proof of Proposition~\ref{prop:packupper} conjugates $H_{\lambda,\alpha,\theta}$ to $H_{\lambda,\alpha,\theta+j\alpha}$ and maps $\delta_j$ to $\delta_0$ and $\{|k-j|>N\}$ to $\{|k|>N\}$; so the bound for the phase $\theta+j\alpha$ is the bound for the start $\delta_j$.
\end{proof}

\subsection{Local dimensions}

A lower bound on transport does not by itself bound packing dimensions from below: the inequalities of Guarneri and Schulz-Baldes go the other way. The lower bound on dimensions comes instead from a pointwise estimate on the resolvent, which the elliptic propagation estimates also provide.

\begin{lemma}[Local mass]\label{lem:localmass}
Let $\theta$, $h$, $c$ and $\eta$ be as in Theorem~\ref{thm:fraction}, and put $h_0=(L+h)/2$. There are $C$ and $n_0$, depending only on $\lambda$, $\alpha$, $h$, $c$ and $\eta$, such that for every $n\ge n_0$ with $b_n\ge c$ and every $E\in U_{h_0}$ with $|\Delta_n(E)|\le2-2\eta$,
\[
|G(0;E+i|\delta_n|)|\le Ce^{h_0q_n},\qquad\mu^\theta\bigl([E-|\delta_n|,E+|\delta_n|]\bigr)\le2Ce^{h_0q_n}|\delta_n| .
\]
\end{lemma}

\begin{proof}
Use Steps~1 and~2 of the proof of Theorem~\ref{thm:fraction} with $T=|\delta_n|^{-1}$, so $\eps=|\delta_n|$ and $\xi=(4\pi\lambda+1)\eps$, on the right half-line only. By~\eqref{eq:poisson} and Lemma~\ref{lem:run}(2),
\[
\operatorname{Im}G(0;z)=\eps\sum_k|G(k;z)|^2\ge\eps\cdot\frac{\eta^{5/2}}{2^{22}}\cdot\frac{\|\Phi_0^+\|^2}{\xi}=c_1\|\Phi_0^+\|^2,\qquad c_1=\frac{\eta^{5/2}}{2^{22}(4\pi\lambda+1)} .
\]
By Lemma~\ref{lem:run}(3), $\operatorname{Im}G(0;z)\le|G(0;z)|\le\|B_0^+\|\,\|\Phi_0^+\|$. Hence $\|\Phi_0^+\|\le\|B_0^+\|/c_1$ and $|G(0;z)|\le\|B_0^+\|^2/c_1\le Ce^{h_0q}$, by the bound on $\|B_0^+\|^2$ in Step~2. Finally $\operatorname{Im}G(0;E+i\eps)\ge\int_{|E'-E|\le\eps}\frac{\eps}{2\eps^2}\,d\mu^\theta(E')$, so $\mu^\theta([E-\eps,E+\eps])\le2\eps\,|G(0;E+i\eps)|$.
\end{proof}

\begin{proof}[Proof of Theorem~\ref{thm:C}, part~1]
Let $\beta=\beta(\alpha)$. Upper local dimensions of a finite measure are nonnegative, so assume $\beta>L$. Fix $h\in(L,\beta)$, $c\in(h,\beta)$ and $n_j\to\infty$ with $b_{n_j}\ge c$ and $b_{n_j}\to\beta$, and put $h_0=(L+h)/2$, $U=U_{h_0}$ and $r_j=|\delta_{n_j}|$.

\emph{Two canonical measures at the same radii.} Let $\theta'=\theta+\alpha$. Translation by one site conjugates $H_{\lambda,\alpha,\theta}$ to $H_{\lambda,\alpha,\theta'}$ and maps $\delta_1$ to $\delta_0$, as in the proof of Proposition~\ref{prop:packupper} below. So $\mu_{\delta_1}=\mu^{\theta'}$, and
\[
\nu:=\mu_{\delta_0}+\mu_{\delta_1}=\mu^\theta+\mu^{\theta'} .
\]
Write $\Delta_n$ and $\Delta_n'$ for the discriminants at the phases $\theta$ and $\theta'$. Since $q_n\theta'=q_n\theta+p_n+\delta_n$, we have $|\Delta_n'-\Delta_n|\le4\pi\lambda^{q_n}|\delta_n|\to0$. For $0<\eta\le1/8$ let $J_j(\eta)=\{E\in U:|\Delta_{n_j}(E)|\le2-4\eta\}$. For large $j$, every $E\in J_j(\eta)$ satisfies $|\Delta_{n_j}(E)|\le2-2\eta$ and $|\Delta'_{n_j}(E)|\le2-2\eta$. Lemma~\ref{lem:localmass}, applied at the phases $\theta$ and $\theta'$, then gives
\[
\nu\bigl([E-r_j,E+r_j]\bigr)\le4Ce^{h_0q_{n_j}}r_j\qquad(E\in J_j(\eta),\ j\text{ large}).
\]

\emph{A full-measure set.} Let $\Lambda_\eta=\limsup_jJ_j(\eta)$. Each $E\in\Lambda_\eta$ lies in $J_j(\eta)$ for infinitely many $j$, and $\log r_j=-b_{n_j}q_{n_j}$, so
\[
\frac{\log\nu([E-r_j,E+r_j])}{\log r_j}\ge1-\frac{h_0}{b_{n_j}}-\frac{\log(4C)}{b_{n_j}q_{n_j}}\to1-\frac{h_0}\beta
\]
along those $j$, with limit $1$ if $\beta=\infty$. Thus $d^+_\nu\ge1-h/\beta$ on $\Lambda_\eta$. For a finite measure, $\nu(\limsup_jJ_j)\ge\limsup_j\nu(J_j)$. The measures $\mu^\theta$ and $\mu^{\theta'}$ are supported on $\Sigma\subset U$, and for large $j$ the set $J_j(\eta)$ contains $\{E\in\Sigma:|\Delta'_{n_j}(E)|\le2-5\eta\}$. Lemma~\ref{lem:arcsine} at the two phases gives
\[
\nu(\Lambda_\eta)\ge\frac2\pi\arcsin(1-2\eta)+\frac2\pi\arcsin\Bigl(1-\frac{5\eta}2\Bigr),
\]
which tends to $2=\nu(\R)$ as $\eta\to0$. So $\bigcup_{k\ge8}\Lambda_{1/k}$ has full $\nu$-measure, and $d^+_\nu\ge1-h/\beta$ holds $\nu$-almost everywhere. Letting $h\downarrow L$ along a sequence gives $d^+_\nu\ge\max\{0,1-L/\beta\}$ $\nu$-almost everywhere.

\emph{Every spectral measure.} Since $\delta_{k+1}=(H-v(\theta+k\alpha))\delta_k-\delta_{k-1}$, the vectors $\delta_0,\delta_1$ generate $\ell^2(\Z)$ under $H$, and every spectral measure $\mu_\varphi$ is absolutely continuous with respect to $\nu$. Let $f=d\mu_\varphi/d\nu$. By differentiation of integrals for Radon measures~\cite[Theorem~2.12 and Corollary~2.14(2)]{Mattila1995}, $\mu_\varphi([E-r,E+r])/\nu([E-r,E+r])\to f(E)$ as $r\downarrow0$ for $\nu$-almost every $E$, and $0<f(E)<\infty$ for $\mu_\varphi$-almost every $E$. At such points $\log\mu_\varphi([E-r,E+r])=\log\nu([E-r,E+r])+O(1)$, so $d^+_{\mu_\varphi}(E)=d^+_\nu(E)$. Since $\mu_\varphi\ll\nu$, the bound $d^+_\nu\ge\max\{0,1-L/\beta\}$ holds $\mu_\varphi$-almost everywhere.
\end{proof}

The upper bound in part~2 is a consequence of the upper transport bound, Theorem~\ref{thm:upper}.

\begin{proposition}\label{prop:packupper}
Let $\lambda>1$ and let $\alpha$ be irrational with $\beta(\alpha)<\infty$. For every $\theta\in\Ga^*$ and every $\varphi\in\ell^2(\Z)$, the spectral measure $\mu_\varphi$ of $H_{\lambda,\alpha,\theta}$ satisfies $\dim_P^+(\mu_\varphi)\le\max\{0,1-\log\lambda/\beta(\alpha)\}$.
\end{proposition}

\begin{proof}
Guarneri and Schulz-Baldes~\cite[Theorem~1]{GSB1999} proved that for every $q>0$ the upper growth exponent of the $q$-th moment, computed with the Ces\`aro average $T^{-1}\int_0^T\cdot\,dt$, is at least the upper packing dimension $\dim_P^+(\mu_{\delta_0})$, defined as the essential supremum of the upper local dimension. The same inequality holds for every unit vector $\psi\in\ell^2(\Z)$ whose spectral measure has compact support, by~\cite[Theorem~1 and Proposition~1(iii)]{BGT2000}: with $\langle\langle|X|^p\rangle\rangle_\psi(T)=T^{-1}\int_0^T\||X|^{p/2}e^{-itH}\psi\|^2\,dt$,
\[
\limsup_{T\to\infty}\frac{\log\langle\langle|X|^p\rangle\rangle_\psi(T)}{\log T}\ge p\,D^+_{\mu_\psi}\Bigl(\frac1{1+p}\Bigr)\ge p\dim_P^+(\mu_\psi)\qquad(p>0),
\]
where $D^+_\mu$ is the upper generalized R\'enyi dimension. Since $(2/T)e^{-2t/T}\ge(2/T)e^{-2}$ for $0\le t\le T$, the Abel-averaged moment $M_q(T)$ dominates $2e^{-2}$ times the Ces\`aro-averaged one. So the Ces\`aro exponents are at most $\beta^+(q)\le\au$, and Theorem~\ref{thm:upper} gives $\dim_P^+(\mu_{\delta_0})\le\max\{0,1-L/\beta\}$ for $\theta\in\Ga^*$.

The translation $(Uu)(n)=u(n+1)$ satisfies $UH_{\lambda,\alpha,\theta}U^{-1}=H_{\lambda,\alpha,\theta+\alpha}$ and $U\delta_1=\delta_0$. Hence the spectral measure of $\delta_1$ for the phase $\theta$ equals that of $\delta_0$ for the phase $\theta+\alpha$, and the same bound holds because $\theta+\alpha\in\Ga^*$ by Lemma~\ref{lem:Gastar}(2). For a Jacobi operator every spectral measure $\mu_\varphi$ is absolutely continuous with respect to $\nu=\mu_{\delta_0}+\mu_{\delta_1}$. The upper packing dimension of a measure is the infimum of the packing dimensions of its full-measure Borel sets~\cite[Appendix]{GSB1999}. Packing dimension of sets is countably stable, so $\dim_P^+(\nu)=\max\{\dim_P^+(\mu_{\delta_0}),\dim_P^+(\mu_{\delta_1})\}$, and absolute continuity gives $\dim_P^+(\mu_\varphi)\le\dim_P^+(\nu)$.
\end{proof}

\begin{proof}[Proof of Theorem~\ref{thm:C}, part~2]
By Proposition~\ref{prop:packupper}, the $\mu_\varphi$-essential supremum of $d^+_{\mu_\varphi}$ is at most $\max\{0,1-L/\beta(\alpha)\}$. By part~1, $d^+_{\mu_\varphi}$ is at least this number $\mu_\varphi$-almost everywhere. So $d^+_{\mu_\varphi}$ equals it $\mu_\varphi$-almost everywhere, and its essential infimum $\dim_P^-(\mu_\varphi)$ and essential supremum $\dim_P^+(\mu_\varphi)$ both equal it. The set $\Ga^*$ has full measure by Lemma~\ref{lem:Gastar}.
\end{proof}
